\documentclass[11pt,aps,preprint,nofootinbib,superscriptaddress,tightenlines]{revtex4-2}
\usepackage[T1]{fontenc}
\usepackage{amsmath,amssymb,mathtools}
\let\lambdabar\relax   % revtex defines it; newpxmath does too
\usepackage{newpxtext,newpxmath}
\usepackage{bm}
\usepackage[final]{microtype}
\usepackage{xcolor}
\usepackage{booktabs}
\usepackage{graphicx}
\usepackage{dsfont}
\usepackage[most]{tcolorbox}
\usepackage{algpseudocode}
\newcommand{\one}{\mathds{1}}
\newcommand{\avg}[1]{\langle #1\rangle}
\definecolor{darkblue}{RGB}{20,60,140}
\definecolor{accent}{RGB}{30,80,150}
\definecolor{tint}{RGB}{243,246,251}
% algorithms: a titled, tinted box
\newtcolorbox{algobox}[1]{enhanced, colback=tint, colframe=accent, boxrule=0.6pt, arc=2pt,
  left=6pt, right=6pt, top=4pt, bottom=4pt, fonttitle=\bfseries\small, coltitle=white,
  colbacktitle=accent, title={#1}, before skip=10pt, after skip=10pt}
% the central equations: a tinted panel with an accent rule on the left
\newtcolorbox{keyeq}{enhanced, colback=tint, frame hidden, borderline west={2pt}{0pt}{accent},
  sharp corners, left=4pt, right=4pt, top=0pt, bottom=0pt, before skip=6pt, after skip=6pt}
\newcommand{\unv}{\textsuperscript{$\dagger$}}
\usepackage[colorlinks=true,linkcolor=darkblue,citecolor=darkblue,urlcolor=darkblue]{hyperref}

\def\st{\begin{equation}}
\def\stp{\end{equation}}
\def\Eq#1{eq.~(\ref{#1})}
\def\Sect#1{Sect.~\ref{#1}}
\def\App#1{Appendix~\ref{#1}}
\def\Fig#1{Fig.~\ref{#1}}
\def\tr{\operatorname{tr}}
\let\oldparagraph\paragraph
\renewcommand{\paragraph}{\par\bigskip\oldparagraph}

\begin{document}

\title{Renormalizing the particle model: the labels and the bare equation of state}
\author{D.~Teaney}
\affiliation{Department of Physics and Astronomy, Stony Brook University,
Stony Brook, New York 11794, USA}
\date{\today}

\begin{abstract}
The particle model conserves the number of particles, which the fluid it represents does not.  Its thermodynamics, computed in \texttt{partition.tex}, is the
input fluid plus an ideal gas of labels and the fluctuations of each particle, of relative order one over the entropy per particle.  This note asks what the
bare (input) equation of state should be.  A static field probes the pressure at fixed $\mu/T$ of the labels; the mechanics uses the pressure as a function of
the energy density, which for the conformal fluid is the input one at every density, so only the temperature scale is renormalized.  No Hamiltonian makes the
pressure at fixed temperature independent of the density of particles.  Number-changing moves, per-particle counterterms and the old entropy volume are
discussed as options.
\end{abstract}

\maketitle
\tableofcontents

%=============================================================================
\section{What the partition function gives}
\label{sec:recap}

From \texttt{partition.tex}, with the notation there ($\bar n$ the density of particles, $\bar V=1/\bar n$, $\sigma=p\bar V/T$, $Z_1$ the partition function of one
particle at fixed volume, $\Pi$ and $\mathcal K$ its pressure and stiffness, $\mathcal K_0=(d+1)/2$).  Only derivatives of $\ln Z_1$ are physical: a constant
factor $\lambda$ per particle in the measure relabels the chemical potential, $\mu/T\to\mu/T-\ln\lambda$, and changes nothing at fixed $N$.

\paragraph{One particle.}  At tree level a particle is a piece of the input fluid, with $f$ its free-energy density as a function of $T$ and the density of particles $n$,
\st
F_1\big|_{\rm tree}=V\,f\Big(T,\frac1V\Big),
\label{Ftree}
\stp
and for the input without a charge $f=-p(T)$.

\paragraph{Pressure and energy.}  With every particle at the mean volume,
\st
p_{\rm model}=\bar nT+p_1(T,\bar V),\qquad E=\bar n\,\avg{e_a}_{\bar V},
\label{pE}
\stp
and at one loop, for the input without a charge, with the fluctuation energy of a particle
\st
\Delta e=d\,T+\frac{T}{c_s^2}\Big[\mathcal K_0-\frac T2\frac{dc_s^2}{dT}\Big]\qquad(=5\,T\ \text{for }d=2,\ c_s^2=\tfrac12),
\label{Deltae}
\stp
\st
p_{\rm model}=p(T)+(1+\mathcal K_0)\,\bar nT,\qquad E=e(T)+\bar n\,\Delta e .
\label{pEloop}
\stp

\paragraph{The labels as a charge.}  The chemical potential of the labels and its one-loop form are
\st
-\frac\mu T=\frac{\partial\ln\mathcal Z_N}{\partial N}\Big|_{T,L}=\ln\bar V+\ln Z_1(T,\bar V)-\Pi,
\label{mu}
\stp
\st
-\frac\mu T=(1+\mathcal K_0)\ln\bar V+\tfrac12\ln\det\big(2\pi\chi(T)\big)-\mathcal K_0,\qquad \chi={\rm diag}\big(C_V,\ T(e+p)\,\delta_{ij}\big),
\label{muloop}
\stp
which for the conformal fluid is a function of the entropy per particle only,
\st
-\frac\mu T=(1+\mathcal K_0)\ln\big(s(T)\bar V\big)+\text{const} .
\label{muconf}
\stp
The model obeys the Euler and Gibbs--Duhem relations of a fluid with one charge,
\st
E+p_{\rm model}=Ts_{\rm model}+\mu\bar n,\qquad dp_{\rm model}=s_{\rm model}\,dT+\bar n\,d\mu ,
\label{GD}
\stp
and the compressibility sum rule gives the long-wavelength fluctuations of the number of particles,
\st
S(0)=T\frac{\partial\bar n}{\partial p_{\rm model}}\Big|_T=\frac{1}{1+\mathcal K} .
\label{S0}
\stp

\paragraph{Hydro-kinetic theory.}  In the hydro-kinetic theory of \cite{amt} the modes below the cutoff renormalize a bare equation of state,
\st
e=e_0+d\,Tn_\Lambda,\qquad p=p_0+\Big[1+\frac T2\frac{dc_s^2}{dT}\Big]Tn_\Lambda ,
\label{amt35}
\stp
and the model has the same three terms with one particle per mode (\App{app:hydrokin}).  As a function of the measured energy density the measured pressure is
\st
p_{\rm model}=p(E)+\Big[1-d\,c_s^2+\frac T2\frac{dc_s^2}{dT}\Big]\bar nT ,
\label{PofE}
\stp
whose bracket vanishes for the conformal fluid.

%=============================================================================

The point: the model's thermodynamics differs from the input at order $1/\sigma$, so the input is a bare equation of state, to be tuned so that the model
reproduces the target.  The model has one more conserved number than the target, the labels, so its equation of state has two variables, and the condition
must say \emph{which} function of the model is the target.  A static field (\Sect{sec:gravity}) and the dynamics (\Sect{sec:offcurve}) answer it: what the
mechanics uses is the pressure as a function of the energy density, which for the conformal fluid is the input one on every curve; the temperature scale is seen
only where the temperature enters explicitly.

\section{What a static field probes}
\label{sec:gravity}

Couple the model to a static potential $\Phi(\bm x)$, slowly varying on the scale $h$:
\st
H\to\sum_ae_a\,\big[1+\Phi(\bm x_a)\big] .
\stp
In equilibrium the temperature redshifts (Tolman), and so does the chemical potential of the labels, which are conserved (Klein):
\st
T(\bm x)\,e^{\Phi(\bm x)}=\text{const},\qquad \frac{\mu(\bm x)}{T(\bm x)}=\text{const} .
\stp
Along this profile the Gibbs--Duhem relation \Eq{GD}, with the Euler relation, gives
\st
dp_{\rm model}=s_{\rm model}\,dT+\bar n\,d\mu=\frac{E+p_{\rm model}}T\,dT ,
\stp
which with $dT/T=-d\Phi$ is the hydrostatic equation $\nabla p=-(E+p)\nabla\Phi$ of a fluid \emph{without} a charge.  So the field sees the model as a
neutral fluid with the equation of state
\begin{keyeq}
\st
P(T)\equiv p_{\rm model}(T,\bar n)\Big|_{\mu/T\ \text{fixed}} .
\label{Pren}
\stp
\end{keyeq}
The bulk modulus of the labels at fixed $T$, $\partial p_{\rm model}/\partial\bar n$, is not probed.

\paragraph{Sound} moves along a different curve in general: the ideal step conserves each $s_a$, so a sound wave holds the entropy per particle fixed.
As in any fluid with a charge, fixed $\mu/T$ and fixed entropy per particle are different curves.  For the conformal fluid they are one, since by
\Eq{muconf} fixed $\mu/T$ is fixed $s(T)\bar V$.

\section{Which entropy per particle}
\label{sec:whichs}

Three entropies per particle appear, and they differ by order one:
\begin{itemize}
\item the mean $\avg{s_a}$ of the particles' own entropy, which the ideal step conserves particle by particle;
\item the bath value $s(T)\bar V=3\sigma$, which the one-loop formulas contain;
\item the start value $S_0\bar V$ of the lattice, by which the runs are labelled.
\end{itemize}
For the conformal fluid the start, which puts all the energy into the rest energy, is hotter by the fluctuation energy, and
\st
S_0\bar V\simeq s(T)\bar V+5
\stp
at one loop; the runs labelled $16$, $54$, $541$ have $s(T)\bar V=11.7$, $49.2$, $536$.  The differences depend on $s(T)\bar V$ alone, so fixing one of the three fixes
the others.  In formulas, $1/s_a$ means $1/(s(T)\bar V)$.

\section{The conformal fluid}
\label{sec:conformal}

Along fixed $\mu/T$ the density of particles follows the entropy, $\bar n\propto s(T)$, and by \Eq{pEloop}
\begin{keyeq}
\st
P(T)=p(T)\Big[1+\frac{1+\mathcal K_0}\sigma\Big],\qquad E(T)=e(T)+5\,\bar nT=2P(T) .
\label{Pconf}
\stp
\end{keyeq}
The model is again conformal, as \Eq{PofE} says.  The labels and the fluctuations renormalize only the normalization of $P$, by $2.5/\sigma$: for a
conformal fluid, the Stefan--Boltzmann constant.  The mechanics does not see it: rescaling the temperature leaves $P(E)=E/2$ unchanged, and in a static field
the energy density follows $E\propto e^{-(d+1)\Phi}$ whatever the normalization.  It is seen where the temperature enters explicitly: a thermometer (the
Tolman profile $T(\bm x)$), and the transport coefficients of the viscous step.

\section{Off the reference curve}
\label{sec:offcurve}

The point: a run does not stay on one curve of fixed $\mu/T$, and no bare equation of state makes $P(T)$ the same on all of them; but the mechanics uses
$P(E)$, which for the conformal fluid is the same on all of them.

\paragraph{Leaving the curve.}  The ideal step conserves each $s_a$, so for the conformal fluid an element that expands or compresses ideally stays on its
curve.  It leaves the curve through the entropy produced by the viscous step, or starts off it if the initial entropy per particle is not uniform.  For a
non-conformal fluid ideal flow crosses the curves as well, since fixed $s_a$ and fixed $\mu/T$ differ.

\paragraph{An exact constraint.}  By Gibbs--Duhem and the compressibility sum rule \Eq{S0}, for any bare Hamiltonian,
\begin{keyeq}
\st
\frac{\partial p_{\rm model}}{\partial\ln\bar n}\Big|_T=\bar nT\,(1+\mathcal K)=\frac{\bar nT}{S(0)} .
\label{constraint}
\stp
\end{keyeq}
$P(T)$ is the same on every curve only if $S(0)=\infty$, the neutral choice of \Sect{sec:notallowed}.  Between two curves at the same temperature, for the conformal fluid
($\bar nT/p=1/\sigma$),
\st
\frac{\Delta P}{P}\Big|_T\simeq-\frac{1+\mathcal K}\sigma\,\Delta\ln s_a ,
\stp
second order in small quantities: $1.6\%$ at $\sigma=16$ for a $10\%$ spread of the entropy per particle.  Lowering $\mathcal K$ with the second dial
(\Sect{sec:dial}) reduces it, at the price of a larger $S(0)$.

\paragraph{What the dynamics uses.}  Ideal hydrodynamics of the model needs the pressure as a function of its conserved densities, the energy and the
labels.  From \Eq{PofE},
\begin{keyeq}
\st
p_{\rm model}(E,\bar n)=p(E)+\Big[1-d\,c_s^2+\frac T2\frac{dc_s^2}{dT}\Big]\bar nT .
\label{PEn}
\stp
\end{keyeq}
For the conformal fluid the bracket vanishes, and $p_{\rm model}=E/2$ at every density: at one loop the labels decouple from the ideal dynamics, and the curves
differ only in the temperature scale, which the mechanics does not see.  The temperature scale enters through the transport coefficients: if $\eta$ is set
through $\eta/s$ with the bare $s(T)$, the model's $\eta/s$ is off by order $(1+\mathcal K_0)/\sigma$ and varies between curves by order
$(1+\mathcal K)\Delta\ln s_a/\sigma$.  The noise needs nothing: the stochastic step is in detailed balance with the model's own ensemble.

\paragraph{Where the labels do couple.}  Two terms make $p_{\rm model}$ depend on $\bar n$ at fixed $E$:
\begin{itemize}
\item the bracket of \Eq{PEn}, the non-conformality of the input times $1/\sigma$;
\item the one-loop term of the positions (\texttt{partition.tex}, ``The positions back in the thermodynamics''), which at fixed $h$ is not conformal.  At $\sigma=16.4$ and $h/a=2$ it gives $\delta p=-0.68\,\bar nT$
  against $\tfrac12\delta E=-0.23\,\bar nT$: a trace anomaly $P-E/2\simeq-0.45\,\bar nT$, about $0.45/\sigma\approx3\%$ of $P$, falling as $1/(\bar nh^2)$.  This is the
  leading coupling of the labels to the mechanics of a conformal model, and the plot of $P/T^3$ along fixed $\mu/T$ (\Sect{sec:measureP}) measures it.
\end{itemize}

\section{The renormalization condition}
\label{sec:dial}

On a reference curve of fixed $\mu/T$, chosen with the resolution, the condition is on the pressure that a static field sees:
\begin{keyeq}
\st
P(T)=p_{\rm model}(T,\bar n)\Big|_{\mu/T\ \text{fixed}}=p_{\rm target}(T) .
\label{rencond}
\stp
\end{keyeq}
Off the curve the error is \Eq{constraint}.  What it fixes depends on the target.
\begin{itemize}
\item \emph{Conformal}: the mechanics is right at one loop on every curve, \Eq{PEn}, with nothing to tune.  \Eq{rencond} is one constant, the temperature scale:
  either the bare pressure is the target divided by $1+(1+\mathcal K_0)/\sigma$, or, equivalently for the dynamics, the transport coefficients are set relative to
  the model's own thermodynamics \Eq{Pconf}.  The constant does not touch the dependence on the volume at fixed $T$, so $\mathcal K$, $S(0)$ and the positions are
  unchanged.
\item \emph{Non-conformal}: \Eq{rencond} is a function of $T$, met by the fixed point below; off the curve the bracket of \Eq{PEn} remains, at order
  $1/\sigma$ times the non-conformality.
\end{itemize}

\paragraph{The general case.}  First choose the resolution: a value of $\mu/T$, equivalently the density $n_0$ at a reference temperature $T_0$.  By
\Eq{muloop}, along the curve the density follows the susceptibilities of the bare equation of state,
\st
n_\mu(T)=n_0\left[\frac{\det\chi(T)}{\det\chi(T_0)}\right]^{1/(d+3)}
\stp
(for the conformal fluid $\det\chi\propto T^{10}$, and $n_\mu\propto T^2\propto s$).  Along it $P=p_{\rm bare}+(1+\mathcal K_0)\,T\,n_\mu$, and \Eq{rencond}
becomes a fixed point,
\begin{keyeq}
\st
p_{\rm bare}(T)=p_{\rm target}(T)-(1+\mathcal K_0)\,T\,n_\mu(T)\big[p_{\rm bare}\big],
\label{pbare}
\stp
\end{keyeq}
with $n_\mu$ computed from the $\chi$ of $p_{\rm bare}$.  The correction is of relative order $1/\sigma$, so one step, with $\chi$ from $p_{\rm target}$,
is right to order $1/\sigma^2$.  The bare function depends on the resolution through $\mu/T$, as $p_0(T;\Lambda)$ in \Eq{amt35} depends on the cutoff.

\begin{algobox}{Dialing the bare equation of state}
\begin{algorithmic}[1]
\State Choose the resolution: the density $n_0$ of particles at a reference temperature $T_0$.
\State Start from the target: $p_{\rm bare}\gets p_{\rm target}$.
\Repeat
  \State From $p_{\rm bare}$ compute the susceptibilities $\chi(T)$ and the density $n_\mu(T)$ along fixed $\mu/T$.
  \State Set $p_{\rm bare}$ to the target less the label and fluctuation pressure along the curve, \Eq{pbare}.
\Until{$p_{\rm bare}$ no longer changes (one or two passes)}
\State At small $\sigma$, use the exact $\ln Z_1$ ((\texttt{partition.tex}, ``Small entropy per particle'')) in \Eq{mu} and \Eq{pE} instead of one loop.
\State Run; measure the mechanical pressure along fixed $\mu/T$ (\Sect{sec:measureP}); correct $p_{\rm bare}$ by the residual $P-p_{\rm target}$, which
  holds the term of the positions (\texttt{partition.tex}, ``The positions back in the thermodynamics''); repeat.
\State Check that $p_{\rm bare}$ is an equation of state: positive entropy and heat capacity, $0<c_s^2<1$.
\end{algorithmic}
\end{algobox}

\paragraph{A second dial.}  Tuning $p_{\rm bare}(T)$ keeps the input without a charge, so the stiffness of the particles keeps its one-loop value and
$S(0)=1/(1+\mathcal K_0)$.  An input with a charge, a free-energy density $f(T,n)$, gives a second dial.  By \Eq{Ftree} its curvature in $n$ adds to the
stiffness at tree level,
\st
\mathcal K=\frac nT\frac{\partial^2f}{\partial n^2}+\mathcal K_0,\qquad S(0)=\frac1{1+\mathcal K},
\stp
so the uniformity of the particles at long wavelength could be chosen, with $P(T)$ along fixed $\mu/T$ still matched by the first dial.  A term in $f$ linear
in $n$ changes only $\mu$, a normalization.  The neutral labels of \Sect{sec:notallowed} are the limit $1+\mathcal K=0$, where the stiffness is left to the gradient term.  Such a term also
changes the forces, through $p_a$.

\section{Holding $\mu/T$ fixed: number-changing moves}
\label{sec:numberchange}

The point: a real fluid is neutral because its number-changing processes are fast, so its chemical potential relaxes instead of being conserved.  The model
could do the same, and then every element would stay on one curve and the reference-curve tuning of \Sect{sec:dial} would hold everywhere.  This is an idea
under consideration, not a plan.

\paragraph{Two versions.}
\begin{enumerate}
\item \emph{Weights moved between neighbours.}  Each particle carries a weight $q_a$, the number of labels it represents, exchanged between neighbours with
  detailed balance, like the entropy in the stochastic step.  This makes $\mu/T$ uniform in space, but the total is conserved, so the common value drifts as
  entropy is produced; and the terms per point, $\bar nT$ and $\mathcal K_0\bar nT$, remain as long as the number of points is fixed.
\item \emph{Particles split and merged} with the grand-canonical weight
  \st
  e^{(\mu/T)\,N}
  \stp
  at a fixed $\mu/T$: the insertion and deletion moves of \Sect{sec:measureP}, now in the dynamics.  Where entropy is produced the entropy per particle rises,
  particles split, and the element returns to the curve.  This is the analogue of number-changing processes, and holds $\mu/T$ at a set value.
\end{enumerate}

\paragraph{What it would buy.}  One curve everywhere, so the dial of \Sect{sec:dial} holds in any flow, conformal or not.  The extra slow mode of the model, the
diffusion of the entropy per particle (\texttt{partition.tex}, Discussion), becomes a relaxing mode.  The labels stay stiff: at fixed $\mu/T$ the number fluctuations are
$\avg{\delta N^2}/N=1/(1+\mathcal K)$.

\paragraph{The Goldstone picture with number-changing moves.}  At a fixed $\mu/T$ the natural object is the grand potential density as a function of the
density of points,
\st
\omega(\bar n)=f(T,\bar n)-\mu\,\bar n,\qquad \omega''(\bar n)=\frac{\partial\mu_{\rm model}}{\partial\bar n}\Big|_T=\frac{T\,\mathcal B}{\bar n},
\stp
with $\mu_{\rm model}$ the chemical potential of the model at the density $\bar n$ and $\mu$ the imposed one.  For neutral labels, $\mathcal B=0$, $\omega$ is
flat: every resolution is a minimum, the vacuum manifold, and a run picks one.  $\mathcal B>0$ curves $\omega$ and selects $\bar n_{\rm eq}(T,\mu/T)$, as the
quark masses tilt the vacuum of QCD.  The relation \Eq{gmor} is then literally ``mass squared times stiffness equals the curvature of the potential at its minimum'',
$m^2f_\rho^2=\omega''$.

The moves change the dynamics, not the statics.  With $N$ conserved the density of points is a conserved order parameter (model B): it relaxes only by moving
points, at a rate proportional to $k^2$.  With number-changing moves it is not conserved (model A) and relaxes in place, at the rate
\st
\Gamma(k)=\Gamma_0\Big(\mathcal B+\frac{\sigma\mu_2k^2}2\Big),\qquad \Gamma(0)=\Gamma_0\,\mathcal B,
\stp
with $\Gamma_0$ the rate of the moves.  The gap in the relaxation is the mass.  So the moves restore the curve only through $\mathcal B$: for neutral labels
there is nothing to relax towards.  Fast number-changing and nearly neutral labels pull against each other, and $\Gamma_0\mathcal B$ is the rate that
matters.

\paragraph{What it would cost.}
\begin{itemize}
\item \emph{The moves must be fast.}  A relaxing internal variable gives a bulk viscosity (Mandelstam--Leontovich), proportional to the relaxation time and to the
  difference between the sound speeds at fixed $s_a$ (frozen) and at fixed $\mu/T$ (equilibrium).  For the conformal fluid the two curves coincide and this vanishes at
  one loop; otherwise it is of order $1/\sigma$, and the moves must relax faster than the hydrodynamic rates of interest.
\item \emph{The moves themselves}: a split or merge must conserve energy and momentum (presumably by adjusting entropies, as the stochastic step does), carry a
  Hastings correction, and handle the change of its neighbours' volumes; it changes the resolution locally, which the fixed-$h$ model does not do now.
\item \emph{The analysis} must follow a fluctuating number of particles.
\end{itemize}

\section{Neutral labels: zero stiffness}
\label{sec:notallowed}

The point: the particles are grid points, and where they are put should not matter.  The neutral choice, a pressure that does not depend on the density of
particles at fixed temperature, says exactly that: at long wavelength moving the points costs nothing.  The points are still held apart on the scale of the
kernel, so the resolution stays controlled, up to a slow wandering.

\paragraph{The condition.}  $p_{\rm model}=p_{\rm target}(T)$ at every density is, by Gibbs--Duhem at fixed $T$ ($dp=\bar n\,d\mu$), the condition $\mu=0$ at
every density.  By \Eq{mu} it fixes the free energy of a particle,
\st
F_1+p_1V=T\ln V\qquad\Longrightarrow\qquad F_1=-pV+T\,(\ln V+1),\qquad \Pi=\sigma-1,\qquad \mathcal K=-1 .
\stp
At one loop this is the counterterm
\st
\delta F_1=T\,(1+\mathcal K_0)\ln V ,
\stp
whose pressure, $-(1+\mathcal K_0)\,T$ per particle, cancels the kinetic stress and the fluctuation share.  It can be put in the Hamiltonian as a shift of the
entropy variable, $e_a(s_a,\bm m_a,V_a)\to e_a\big(s_a+(1+\mathcal K_0)\ln V_a,\bm m_a,V_a\big)$, which shifts $\ln Z_1$ by $-(1+\mathcal K_0)\ln V$ exactly; it
changes the forces.  The labels are then neutral: $\mu/T$ is the same on every curve, $P(T)$ is the same on every curve, the bare equation of state is the physical
one up to the one-loop term of the positions, and entropy production no longer moves a run off its curve (\Sect{sec:offcurve}).

\paragraph{The positions.}  With $\mathcal K=-1$ the kernel of the structure factor, $S=1/(1+A)$ with $A=2\Pi\widetilde W(1-\widetilde W)+\mathcal K\widetilde W^2$, factorizes,
\st
1+A(k)=\big(1-\widetilde W\big)\big[1+(1+2\Pi)\,\widetilde W\big],
\stp
positive at every $k$ since $0\le\widetilde W<1$: every mode is stable.  At $kh\ll1$, with $\widetilde W\simeq1-\mu_2k^2/4$,
\begin{keyeq}
\st
S(k)\simeq\frac{2}{\sigma\,\mu_2k^2} .
\stp
\end{keyeq}
The points are held at long wavelength not by counting but by the fluid's pressure acting through the kernel, a gradient stiffness.  This is the middle regime
that the model already has between $h$ and the label-correlation length $\xi$ (\texttt{partition.tex}); the counterterm sends $\xi$ to infinity.  The run at
$\sigma=179$, with $\xi=4.5h$ in a box of $9h$, is already mostly in this regime.  The points behave like an adaptive mesh with a smoothness penalty, free to
follow the flow at long wavelength.

\paragraph{The price: the resolution wanders, slowly.}  With $S\propto1/k^2$ the relative fluctuation of the number of points in a region of size $R$, in a box of
size $L$, is
\st
\frac{\avg{\delta N_R^2}}{N_R^2}\simeq\frac{\ln(L/R)}{\pi\,\sigma\,\bar n\mu_2} ,
\stp
independent of $R$ apart from the logarithm: about $10\%$ at $\sigma=16$ and $h/a=2$, and $3\%$ at $\sigma=179$.  The volume variance grows in the same way,
as the logarithm of the box in $d=2$; in $d=3$ both are finite.  So the number of neighbours varies at the level of $(\sigma\bar nh^2)^{-1/2}$, with a logarithm.

\paragraph{A regulator.}  The stiffness need not be removed completely.  Counter only most of the label term, leaving a small residual stiffness
\st
\mathcal B\equiv1+\mathcal K\ll1,\qquad \delta F_1=T\,(1+\mathcal K_0-\mathcal B)\ln V .
\stp
The label pressure is then the regulator alone,
\st
p_{\rm model}=p(T)+\mathcal B\,\bar nT ,
\stp
and the structure factor is of the Ornstein--Zernike form, with a length set by the regulator,
\begin{keyeq}
\st
1+A=\big(1-\widetilde W\big)\big[1+(1+2\Pi)\,\widetilde W\big]+\mathcal B\widetilde W^2,\qquad
S(k)\simeq\frac{1/\mathcal B}{1+\xi^2k^2},\qquad \xi^2=\frac{\sigma\mu_2}{2\mathcal B} .
\label{regulator}
\stp
\end{keyeq}
Both costs scale with $\mathcal B$: the logarithm in the wandering of the resolution is cut off at $\xi$ instead of the box, $\ln(L/R)\to\ln(\xi/R)$,
and the mismatch between curves is
\st
\frac{\Delta P}{P}\Big|_T\simeq-\frac{\mathcal B}\sigma\,\Delta\ln s_a .
\stp
So $\mathcal B$ is a regulator: choose $\xi$, say a fraction of the box, and $\mathcal B$ follows.  The present model has
$\mathcal B=1+\mathcal K_0=5/2$; the neutral labels are $\mathcal B\to0$.

\paragraph{The symmetry and its Goldstone mode.}  $S\propto1/k^2$ is the correlator of a massless field, as for the phase of a condensate.  Here the
massless field is the density of grid points itself, and the symmetry is the freedom to choose the resolution.

\emph{The symmetry.}  At fixed temperature and fixed fluid fields, change the density of grid points by a constant,
\st
\rho(\bm x)\to\rho(\bm x)+c .
\stp
A uniform change of the density of points changes the weight $\Phi[\rho]$ (\texttt{partition.tex}) by the free energy of the labels, whose density-dependent
part at fixed $T$ is set by the bulk modulus of the labels: from \Eq{pE} and \Eq{GD},
\st
\frac{\partial(\mu/T)}{\partial\bar n}\Big|_T=\frac{1+\mathcal K}{\bar n}=\frac{\mathcal B}{\bar n} .
\stp
For neutral labels, $\mathcal B=0$, the chemical potential does not depend on the density, the free energy is linear in the number of points, and the
shift changes $\Phi$ only by $(\mu/T)\,\delta N$, a constant at fixed $N$: every resolution is equally weighted.  This is the compressional half of the
relabelling symmetry of a fluid, the freedom to move the grid points without changing the fluid.  The other half, the moves at fixed local density of points
(transverse displacements, $\nabla\!\cdot\bm u=0$), is a symmetry of the model at any $\mathcal K$: in the Gaussian theory they cost nothing.  With neutral
labels the whole group of smooth relabellings, compressional included, is a symmetry at long wavelength.

\emph{The Goldstone mode.}  The resolution is then a modulus: the run picks a value, $N/L^2$, but nothing in the weight prefers it locally, so different
regions may sit at different resolutions.  Its slow variations are the Goldstone mode.  They are not free, because the fluid's pressure sees a difference
of resolution through the kernel, the term $\widetilde W(1-\widetilde W)\simeq\mu_2k^2/4$ of $A(k)$, which is a gradient.  At long wavelength
\begin{keyeq}
\st
\Phi[\rho]\simeq\text{const}-\frac1{2\bar n}\int d^2x\,\Big[\mathcal B\,\delta\rho^2+\frac{\sigma\mu_2}2\,(\nabla\delta\rho)^2\Big],
\qquad S(k)=\frac{1}{\mathcal B+\sigma\mu_2k^2/2} ,
\label{goldstone}
\stp
\end{keyeq}
the action of a scalar with the stiffness $\sigma\mu_2/2$ and the mass squared $\mathcal B$.  The shift symmetry forbids the mass term; the gradient term
is allowed, and is what remains.

\emph{What the gradient costs, and what it is not.}  The gradient term is the fluid's pressure acting on the volume that the kernel fails to assign:
across a gradient of resolution the volumes of the particles do not tile the box, and to second order
\st
\sum_aV_a-L^2=-\frac1{\bar n^2L^2}\sum_{\bm k}\widetilde W\big(1-\widetilde W\big)\,|\delta\rho_{\bm k}|^2,\qquad \widetilde W\big(1-\widetilde W\big)\simeq\frac{\mu_2k^2}4 ,
\stp
weighted by $p/T$.  It is a stiffness of the density, not an elastic rigidity.  In terms of the displacement of the points, $\delta\rho=-\bar n\,\nabla\!\cdot\bm u$,
\st
\frac{\sigma\mu_2}{4\bar n}\,(\nabla\delta\rho)^2=\frac{\sigma\mu_2\bar n}4\,\big(\nabla\nabla\!\cdot\bm u\big)^2\;\propto\;k^4\,|u_{L,\bm k}|^2 ,
\stp
a bending-type term, two powers of $k$ softer than the $k^2|\bm u_{\bm k}|^2$ of an elastic solid; transverse displacements cost nothing.  So in the Gaussian
theory the points have neither a shear nor a longitudinal elastic modulus.  Any rigidity of the arrangement (the open question of a transverse sound speed that
falls as $h/a$ grows) comes from beyond the Gaussian theory, from the core of $g(r)$ near $k_*$ where the points order locally, and should fade as the number
of neighbours grows.

\emph{The regulator is the mass.}  $\mathcal B$ breaks the symmetry explicitly, as a small field does in the XY model, and $\xi$ of \Eq{regulator} is the
inverse mass.  The present model breaks it strongly, $\mathcal B=1+\mathcal K_0=5/2$: the counting of points (the $1$) and the stiffness of each particle
(the $\mathcal K_0$) pin the resolution to its mean.

\emph{A Gell-Mann--Oakes--Renner relation.}  In QCD the pion mass, the pion decay constant and the explicit breaking of chiral symmetry are tied by
\st
m_\pi^2f_\pi^2=(m_u+m_d)\,|\avg{\bar qq}| :
\stp
the mass squared of the Goldstone boson times its stiffness equals the curvature of the vacuum energy along the symmetry direction, which is linear in the
explicit breaking.  The same three quantities appear here.  Write \Eq{goldstone} in units of $T$, so that it is a free energy:
\begin{itemize}
\item \emph{the stiffness} (the analogue of $f_\pi^2$) is the coefficient of the gradient term,
  \st
  f_\rho^2\equiv\frac{T\,\sigma\mu_2}{2\bar n}=\frac{p\,\mu_2}{2\bar n^2} ,
  \stp
  set by the fluid's pressure and the size of the kernel, and independent of the breaking;
\item \emph{the curvature along the symmetry direction} (the analogue of $(m_u+m_d)|\avg{\bar qq}|$) is the cost of a uniform change of the density of points, the
  inverse compressibility of the labels,
  \st
  \frac{\partial\mu}{\partial\bar n}\Big|_T=\frac{T\,\mathcal B}{\bar n}=\frac{T}{\bar n\,S(0)} ,
  \stp
  linear in the explicit breaking $\mathcal B$;
\item \emph{the mass} is $m=1/\xi$.
\end{itemize}
Then \Eq{regulator} reads
\begin{keyeq}
\st
m^2f_\rho^2=\frac{\partial\mu}{\partial\bar n}\Big|_T=\frac{T\,\mathcal B}{\bar n} .
\label{gmor}
\stp
\end{keyeq}
As for the pion, the mass squared is linear in the breaking, $m^2\propto\mathcal B$, not its square root, and the stiffness does not depend on it at
leading order.  The relation that guarantees it is the analogue of the Ward identity behind the Gell-Mann--Oakes--Renner relation: the compressibility sum
rule \Eq{S0}, which ties the $k\to0$ limit of the correlator to the thermodynamic curvature exactly.  The present model, $\mathcal B=5/2$, is the
analogue of heavy quarks: the ``Goldstone'' mass $1/\xi$ is of the order of the inverse kernel size at small $\sigma$ ($\xi=0.66h$ at $\sigma=3.9$), and the
mode is not light.  The regulator is the analogue of light quarks.

\emph{As in two dimensions.}  Without the mass the fluctuations grow as the logarithm of the box (the wandering of the resolution above), as for a condensate in
$d=2$ (Mermin--Wagner): only quasi-long-range order of the density of points.  The displacements of the points are softer still: with
$\delta\rho_{\bm k}=-i\bar n\,\bm k\cdot\bm u_{\bm k}$, $\avg{|u_{L,\bm k}|^2}\propto1/k^4$, as for a membrane rather than a phonon.

\emph{Where the analogy stops.}  In a superfluid the massless field is the phase and the density is gapped, and the symmetry is a physical one, with a conserved
current and second sound.  Here the massless field is the density of grid points, and the symmetry is a redundancy of the description, the choice of the grid.
Nothing physical is attached to the Goldstone mode: the fluid fields see it only through the kernel, at the order of the gradient term.

\paragraph{Open.}  The terms beyond the Gaussian at small $\mathcal B$, which matter at long wavelength once the quadratic stiffness is small, are not
checked.  The exact (all orders in $1/s_a$) version needs $\ln Z_1$ from quadrature, since the shift $(1+\mathcal K_0)\ln V$ removes the stiffness at one loop only.
Shear rigidity, a transverse sound speed of the arrangement, is a separate question: $S(0)$ is the bulk response of the labels, and in the Gaussian theory
transverse displacements of the points cost nothing at any $\mathcal K$.

\paragraph{The value of $\mu/T$.}  A fixed value of $\mu/T$, zero included, only labels a curve, by the normalization of the measure.  What the neutral choice
fixes is not the value but its independence of the density.

\section{Measuring the pressure at fixed $\mu/T$}
\label{sec:measureP}

Two routes.  With a common Lagrange multiplier, the grand-canonical weight $e^{(\mu/T)N}$ holds $\mu/T$ fixed directly; this needs moves that insert and delete
particles.  Without new moves, canonical runs on a grid of temperatures and densities give $\mu/T$ up to one constant through
\st
\frac{\partial(\mu/T)}{\partial\bar n}\Big|_T=\frac1{\bar nT}\frac{\partial p_{\rm model}}{\partial\bar n}\Big|_T,\qquad
\frac{\partial(\mu/T)}{\partial T}\Big|_{\bar n}=-\frac1{T^2}\frac{\partial E}{\partial\bar n}\Big|_T ,
\label{maxwell}
\stp
the first from Gibbs--Duhem, the second a Maxwell relation of $\ln\mathcal Z_N$.  The missing constant is the relabelling of the curves, and is harmless.

\begin{algobox}{The pressure at fixed $\mu/T$}
\begin{algorithmic}[1]
\If{grand canonical}
  \State Add to the hybrid Monte Carlo moves that insert and delete a particle with the weight $e^{\mu/T}$ per particle, the new particle's entropy and
  momentum drawn from the single-particle distribution at its trial volume; the change of weight is local, through the volumes of its neighbours.
  \State For each $T$: run, and record the mean density and the mechanical pressure ($\tau^{\rm kin}$ plus the virial).
\Else
  \State For each $T$ and $\bar n$ on a grid: run canonically; record the mechanical pressure and the energy density.
  \State Integrate \Eq{maxwell} for $\mu/T$ on the grid, up to one constant.
  \State Interpolate the pressure along the contours of $\mu/T$.
\EndIf
\State Plot $P/T^3$ against $T$ on each curve: constant for an exactly conformal model.
\end{algorithmic}
\end{algobox}

%=============================================================================
%=============================================================================
\appendix

\section{Per-particle counterterms and the entropy volume}
\label{app:entropyvolume}

The point: the cutoff terms of the model are of order $T$ per particle, so the natural counterterms are per particle, terms in $H$ of the form $-\mu N$,
rather than changes of the energy density per unit volume.  This appendix records what such terms can and cannot do, and the old entropy volume, which applies
the same idea to the volumes.  Notes for later; nothing here goes beyond one loop.

\subsection{The bare energy and the number of particles}

In the bare Hamiltonian of the n-scheme the energy density does not depend on the density of particles at fixed entropy and momentum (\texttt{volume\_n.tex}),
\st
\mu_n=\Big(\frac{\partial h}{\partial n}\Big)_{S,\bm M}=0 .
\stp
This is a property of the bare Hamiltonian only.  The model's own thermodynamics has $\mu\neq0$ at order $T$ per particle, \Eq{muloop}, and its pressure
depends on the density at fixed $T$, \Eq{constraint}.  The counterterms that remove cutoff terms of order $T$ per particle are therefore naturally per particle,
\st
H\to H+\sum_a\delta\varepsilon_a .
\stp

\paragraph{A constant per particle,} $\delta\varepsilon_a=-\mu_c$, is the term $H-\mu_cN$.  At fixed $N$ it changes nothing but the normalization: it shifts the
chemical potential of the labels by $-\mu_c$, which relabels the curves (\Sect{sec:recap}).  It matters as soon as $N$ changes.  With number-changing moves
(\Sect{sec:numberchange}) the moves see $H-\mu N$, and $\mu_c$ is the dial that sets the resolution, the curve on which the run sits.

\paragraph{A term that depends on the particle's state,} $\delta\varepsilon_a=\delta\varepsilon(s_a,V_a)$ of order $T$, changes $\ln Z_1$ by $-\avg{\delta\varepsilon}_V/T$
at leading order, and so $\Pi$, $\mathcal K$ and $\Delta e$ by numbers of order one per particle, like $\mathcal K_0$.  Two things it cannot do:
\begin{itemize}
\item \emph{evade the constraint} \Eq{constraint}, which holds for any Hamiltonian: $P(T)$ differs between curves unless $S(0)=\infty$;
\item \emph{cancel the temperature scale on every curve} while keeping the model conformal.  At one loop the label and fluctuation terms of $E$ and of $p_{\rm model}$,
  \st
  \delta E=5\,\bar nT,\qquad \delta p_{\rm model}=\tfrac52\,\bar nT ,
  \stp
  cancel in $P-E/2$ only together; removing the energy alone breaks $P=E/2$.
\end{itemize}
What it can do is set the stiffness $\mathcal K$: the second dial of \Sect{sec:dial}, whose terms in $f(T,n)$ beyond the linear one are a free energy per
particle that depends on the density.

\subsection{The entropy volume}

Before 2026-11-10 the volume was the particle's share of the smoothed entropy density (\texttt{SPH\_VOLUME=entropy}, \texttt{volume\_n.tex}),
\st
V_a=\frac{s_a}{S_a},\qquad S_a=\sum_bs_b\,W_{ab} .
\stp
It weights each particle's share by its own entropy, a physical variable that the viscous step does not conserve.  (For any weights $q_a$ in place of $s_a$ only the overall scale
is a gauge: rescaling all of them by one factor leaves every volume unchanged.  Their ratios are physical.)

\paragraph{What it gets right.}  In the n-scheme the entropy per particle is a passive charge with almost no restoring force, and a particle with more entropy
must push its neighbours away to reach the common pressure: the spread of $s_a$ becomes disorder of the positions (\texttt{volume\_n.tex}).  With the entropy
volume the weights make the same adjustment at once, $V_a\propto s_a$, and the positions can stay ordered.  The spread that the viscous step produces is absorbed
by the volumes, not by the arrangement.

\paragraph{What it does not fix.}  The points are still conserved, and each carries the kinetic stress $d\,T$ and a fluctuation share of order $T$.  A particle
whose weight, and so whose volume, goes to zero does not drop out: equipartition, $\avg{\bm m_a\cdot\bm v_a}=d\,T$, holds for any energy function, and as
$V_a\to0$ at fixed momentum the momentum density diverges, the particle becomes ultrarelativistic, $e_a\simeq|\bm m_a|$, and
\st
\avg{|\bm m_a|}=d\,T :
\stp
a massless quantum with energy and kinetic stress of order $T$.  The gas of labels is, in this sense, a gas of such quanta, one per point.  Such particles
are rare, since the weight of a particle carries $V_a^{\mathcal K_0}$; removing them is a number-changing move.  Elements with
different entropy per point still sit on different curves of fixed $\mu/T$, with the mismatch of \Sect{sec:offcurve}.

\paragraph{Why it was dropped.}  The stochastic step changes the weights at fixed volume, but the volumes depend on the weights, so the work of their instant
redistribution,
\st
-\Delta p\,\Delta V ,
\stp
was lost at every step.

\paragraph{Open questions.}
\begin{enumerate}
\item \emph{Its partition function.}  At fixed positions the particles are not independent, since each volume depends on its neighbours' entropies, so $\mathcal Z_N$
  does not factorize into $\prod_aZ_1(V_a)$ and the theory of this note does not carry over.  A one-loop calculation would give what replaces $\mathcal K_0$, the
  stiffness of the points and $S(0)$, and the label pressure.
\item \emph{Its back map.}  Whether an energy-conserving version, with the volumes moved inside each move and the back map solved at the new volumes, fits the
  inverse-moment step and its Hastings correction.
\end{enumerate}

\section{The hydro-kinetic comparison in detail}
\label{app:hydrokin}

Eq.~(35) of \cite{amt} is, with $n_\Lambda=\Lambda^3/6\pi^2$ the number of modes per volume,
\st
e(T_0)=e_0(T_0;\Lambda)+3\,Tn_\Lambda,\qquad p(T_0)=p_0(T_0;\Lambda)+\Big[1+\frac T2\frac{dc_s^2}{dT}\Big]Tn_\Lambda .
\label{amt35}
\stp
One mode corresponds to one particle, $n_\Lambda\leftrightarrow\bar n$, and the three terms are the same.

\begin{table}[h]
\caption{The terms of \Eq{amt35} in the particle model.}
\begin{ruledtabular}
\begin{tabular}{ll}
hydro-kinetic theory & SPH model \\
\colrule
$\avg{G^2}/(e+p)=d\,Tn_\Lambda$ & kinetic energy of $\bm m_a$, $d\,T\bar n$ \\
$\avg{G^iG^i}/d(e+p)=Tn_\Lambda$ & kinetic stress $\avg{m^iv^j}$, the $1$ in $1+\mathcal K_0$ \\
$\tfrac12p''(e)\avg{\delta e^2}=\tfrac T2\tfrac{dc_s^2}{dT}\,Tn_\Lambda$ & $\avg{p_a}-p_{\rm in}(\bar e_{\rm rest})$\unv \\
\end{tabular}
\end{ruledtabular}
\end{table}

\paragraph{The quantities.}  Four energies and three temperatures appear, and the notations of \cite{amt} and of \texttt{partition.tex} clash
(there $T_0$ is the bath temperature of the runs).  Here:
\begin{itemize}
\item $T$ is the \emph{bath temperature}: the parameter of the canonical ensemble.  It is the thermodynamic temperature of the model, the one in
  $\avg{\delta E^2}=T^2\partial E/\partial T$.
\item $e_{\rm in}(T)$, $p_{\rm in}(T)$ are the \emph{input} equation of state, the functions in the code, defined through $T_a=\partial e_a/\partial s_a$.  They are
  the $e$ and $p$ of the rest of this note.
\item $e_{{\rm rest},a}$ is the rest-frame energy density of a particle, a function of $(s_a,\bm m_a,V_a)$ through the input equation of state.
  To second order in the momentum the energy density on the slice is $e_{{\rm rest},a}+\bm M_a^2/(e+p)$.
\item $\bar e_{\rm rest}(T)\equiv\avg{e_{{\rm rest},a}}$ is its mean in the ensemble at $T$.  It is the $e_0$ of \cite{amt}: the mean energy density with the kinetic
  energy of the fluctuations removed.
\item $e_{\rm tot}(T)=\bar e_{\rm rest}+d\,T\bar n$ is the measured energy density.
\end{itemize}
The mean is not the input function at the bath temperature.  From $e_{\rm tot}=e_{\rm in}(T)+\bar n\Delta e$ and \Eq{Deltae},
\begin{keyeq}
\st
\bar e_{\rm rest}(T)=e_{\rm in}(T)+\bar n\,\frac{T}{c_s^2}\Big[\mathcal K_0-\frac T2\frac{dc_s^2}{dT}\Big] .
\label{erest}
\stp
\end{keyeq}
This is the difference between the mean and the most probable value of a skewed distribution.  For the energy alone, with a flat measure in
$e$, it is the textbook $\kappa_3/2\kappa_2=T+T^2C_V'/2C_V$; the $\mathcal K_0$ here collects one half from the entropy and $d/2$ from the momenta, whose
Gaussian width depends on the rest energy.

\paragraph{Two labels.}  A temperature can be assigned to an energy density through the input function.  Two such labels are in use:
\st
e_{\rm in}(T_e)\equiv\bar e_{\rm rest},\qquad e_{\rm in}(T_{\rm start})\equiv e_{\rm tot},\qquad\qquad T<T_e<T_{\rm start},
\stp
\st
T_e-T=\frac T{s_a}\Big[\mathcal K_0-\frac T2\frac{dc_s^2}{dT}\Big]\unv,\qquad
T_{\rm start}-T=\frac{c_s^2\,\Delta e}{s_a}=\frac T{s_a}\Big[d\,c_s^2+\mathcal K_0-\frac T2\frac{dc_s^2}{dT}\Big] .
\label{Tshift}
\stp
$T_e$ and $T_{\rm start}$ are labels of energies, not temperatures.  $T_{\rm start}$ is the $T(S_0)$ of the lattice start in \texttt{partition.tex}, where all of the
energy is rest energy; the bath temperature of the runs is below it by the second of \Eq{Tshift}.

\paragraph{The comparison.}  In \cite{amt} the bare function is \emph{defined} by the mean, $e_0(T;\Lambda)\equiv\bar e_{\rm rest}(T)$, so that there the label $T_e$ and the
temperature coincide and are both called $T_0$.  In the particle model the bare function is the input, and they differ.  The content is the
same.  Written with the mean rest energy, \Eq{amt35} holds as it stands,
\st
e_{\rm tot}=\bar e_{\rm rest}+d\,T\bar n,\qquad p_{\rm mech}=p_{\rm in}(\bar e_{\rm rest})+\Big[1+\frac T2\frac{dc_s^2}{dT}\Big]T\bar n ,
\stp
and contains no $\mathcal K_0$: the relation between the measured pressure and the measured energy does not depend on the measure.  Written with
the input functions at the bath temperature, using \Eq{erest} and $dp_{\rm in}=c_s^2\,de_{\rm in}$,
\st
e_{\rm tot}=e_{\rm in}(T)+\bar n\,\Delta e,\qquad p_{\rm mech}=p_{\rm in}(T)+(1+\mathcal K_0)\,\bar nT ,
\stp
with \Eq{Deltae}, which reproduces $5\,T$.  $\mathcal K_0$ enters only the assignment of a temperature to an energy; the derivative of $c_s^2$ sits in the
pressure at fixed mean energy and in the energy at fixed bath temperature.  The label terms have no counterpart in \cite{amt}.

\paragraph{The name of $\mathcal K_0$.}  $\mathcal K_0=(d+1)/2$ is half the number of fluctuating variables integrated per particle.  It is not a compressibility of the
fluid.  It is tied to that of the labels, $\bar nT\kappa_T^{\rm label}=S(0)=1/(1+\mathcal K_0)$.

\begin{thebibliography}{9}
\bibitem{amt} Y.~Akamatsu, A.~Mazeliauskas and D.~Teaney, \emph{Bulk viscosity from hydrodynamic fluctuations with relativistic hydro-kinetic theory},
  Phys.\ Rev.\ C \textbf{97}, 024902 (2018), arXiv:1708.05657.
\end{thebibliography}

\end{document}
